% claim.tex -- cyl3-mod23 priority claim note (SOP 08b claim channel).
% AI-generated 2026-09-27 by the AI4Math paper department (claim lane).
% lstlean.tex copied unchanged from harness/templates/paper/lstlean.tex.
% Machine-readable anchors: every theorem carries a % LEAN: line (task id +
% declaration name + grade); paper-lint.sh and the claim audit check against it.
% All Lean listings are byte-exact mechanical splices of the frozen artifacts
% of tasks/20260926-cyl3-pipeline (see audit/listing-verify.txt in this
% package). The listing's Chinese header comments are pipeline provenance
% notes (not mathematical content); they render through the package's CJK
% fontset (xeCJK/SimSun), a rendering-only preamble addition copied from the
% sanctioned claims/tatami-mod8-defect pattern and extended to this file's
% symbol inventory. Frozen source bytes are unchanged.
\documentclass[11pt]{article}

\usepackage[T1]{fontenc}
\usepackage[utf8]{inputenc}
\usepackage{amsmath,amssymb,amsthm}
\usepackage{hyperref}
\usepackage{graphicx}
% lstlean.tex — Lean style for the listings package.
% Lineage: Jeremy Avigad (2015), adapted from Assia Mahboubi's SSR style;
% inspired by Olivier Verdier's unixode.sty for the Unicode literate table.
% No CTAN package or upstream repo exists; papers carry this file in their
% source package (cap set arXiv:1907.01449, sphere eversion arXiv:2210.07746).
% AI4Math adaptation (AI-generated, 2026-09-26): sorry/admit/sorryAx elevated
% to a dedicated error tier, matching the pipeline's 0-sorry constitution.
\usepackage{listings}
\usepackage{xcolor}

\definecolor{keywordcolor}{rgb}{0.7, 0.1, 0.1}   % red keywords
\definecolor{tacticcolor}{rgb}{0.0, 0.1, 0.6}    % blue tactics
\definecolor{commentcolor}{rgb}{0.4, 0.4, 0.4}   % grey comments
\definecolor{symbolcolor}{rgb}{0.0, 0.1, 0.6}    % blue symbols
\definecolor{sortcolor}{rgb}{0.1, 0.5, 0.1}      % green sorts
\definecolor{errorcolor}{rgb}{0.6, 0, 0}         % dark red: sorry/admit
\definecolor{stringcolor}{rgb}{0.5, 0.1, 0.5}    % purple strings

\lstdefinelanguage{lean}{
  % Anything between $ becomes LaTeX math mode
  mathescape=false,
  % Comments may or not include Latex commands
  texcl=false,
  % keywords, list taken from lean-syntax.txt
  morekeywords=[1]{import, prelude, protected, private, noncomputable,
    definition, meta, renaming, hiding, parameter, parameters, begin,
    constant, constants, lemma, variable, variables, theory, print,
    theorem, example, open, section, namespace, universe, universes, end,
    modifier, modifiers, using, export, exposing, axiom, inductive,
    coinductive, with, structure, class, instance, attribute, local, where,
    by, have, show, let, in, if, then, else, match, do, for, fun, assume,
    from, take, calc, include, omit, infix, infixl, infixr, notation,
    macro, syntax, elab, set_option, initialize, deriving, opaque, abbrev,
    mutual, partial, scoped, termination_by, decreasing_by},
  % Sorts
  morekeywords=[2]{Sort, Type, Prop},
  % tactics, list taken from lean-syntax.txt (tactic keywords)
  morekeywords=[3]{intro, intros, apply, exact, refine, rfl, simp, simp_all,
    simpa, simp_rw, rw, rewrite, erewrite, ring, ring_nf, omega, decide,
    norm_num, norm_cast, induction, cases, rcases, obtain, constructor,
    ext, funext, conv, congr, field_simp, linarith, nlinarith, positivity,
    tauto, trivial, aesop, use, by_cases, by_contra, push_neg, contrapose,
    symm, trans, assumption, contradiction, exfalso, specialize,
    generalize, revert, clear, rename_i, subst, unfold, delta, change,
    convert, split_ifs, interval_cases, fin_cases, zify, gcongr,
    nth_rewrite, suffices, generalizing, at, only, native_decide},
  % warnings - constitution tier: must never survive into a final proof
  morekeywords=[4]{sorry, admit, sorryAx},
  literate=
    {α}{{\ensuremath{\alpha}}}1
    {β}{{\ensuremath{\beta}}}1
    {γ}{{\ensuremath{\gamma}}}1
    {δ}{{\ensuremath{\delta}}}1
    {ε}{{\ensuremath{\varepsilon}}}1
    {ζ}{{\ensuremath{\zeta}}}1
    {η}{{\ensuremath{\eta}}}1
    {θ}{{\ensuremath{\theta}}}1
    {ι}{{\ensuremath{\iota}}}1
    {κ}{{\ensuremath{\kappa}}}1
    {λ}{{\ensuremath{\lambda}}}1
    {μ}{{\ensuremath{\mu}}}1
    {ν}{{\ensuremath{\nu}}}1
    {ξ}{{\ensuremath{\xi}}}1
    {π}{{\ensuremath{\pi}}}1
    {ρ}{{\ensuremath{\rho}}}1
    {σ}{{\ensuremath{\sigma}}}1
    {τ}{{\ensuremath{\tau}}}1
    {υ}{{\ensuremath{\upsilon}}}1
    {φ}{{\ensuremath{\varphi}}}1
    {χ}{{\ensuremath{\chi}}}1
    {ψ}{{\ensuremath{\psi}}}1
    {ω}{{\ensuremath{\omega}}}1
    {Γ}{{\ensuremath{\Gamma}}}1
    {Δ}{{\ensuremath{\Delta}}}1
    {Θ}{{\ensuremath{\Theta}}}1
    {Λ}{{\ensuremath{\Lambda}}}1
    {Ξ}{{\ensuremath{\Xi}}}1
    {Π}{{\ensuremath{\Pi}}}1
    {Σ}{{\ensuremath{\Sigma}}}1
    {Φ}{{\ensuremath{\Phi}}}1
    {Ψ}{{\ensuremath{\Psi}}}1
    {Ω}{{\ensuremath{\Omega}}}1
    {ℵ}{{\ensuremath{\aleph}}}1
    {∀}{{\ensuremath{\forall}}}1
    {∃!}{{\ensuremath{\exists}!}}1
    {∃}{{\ensuremath{\exists}}}1
    {∈}{{\ensuremath{\in}}}1
    {∉}{{\ensuremath{\notin}}}1
    {∘}{{\ensuremath{\circ}}}1
    {∩}{{\ensuremath{\cap}}}1
    {∪}{{\ensuremath{\cup}}}1
    {⊆}{{\ensuremath{\subseteq}}}1
    {⊂}{{\ensuremath{\subset}}}1
    {⊇}{{\ensuremath{\supseteq}}}1
    {⊃}{{\ensuremath{\supset}}}1
    {⊄}{{\ensuremath{\nsubseteq}}}1
    {∅}{{\ensuremath{\emptyset}}}1
    {∧}{{\ensuremath{\wedge}}}1
    {∨}{{\ensuremath{\vee}}}1
    {¬}{{\ensuremath{\neg}}}1
    {⊤}{{\ensuremath{\top}}}1
    {⊥}{{\ensuremath{\bot}}}1
    {↔}{{\ensuremath{\leftrightarrow}}}1
    {→}{{\ensuremath{\rightarrow}}}1
    {←}{{\ensuremath{\leftarrow}}}1
    {↑}{{\ensuremath{\uparrow}}}1
    {⇒}{{\ensuremath{\Rightarrow}}}1
    {⇐}{{\ensuremath{\Leftarrow}}}1
    {⟹}{{\ensuremath{\Longrightarrow}}}1
    {↦}{{\ensuremath{\mapsto}}}1
    {≤}{{\ensuremath{\leq}}}1
    {≥}{{\ensuremath{\geq}}}1
    {≠}{{\ensuremath{\neq}}}1
    {≈}{{\ensuremath{\approx}}}1
    {≡}{{\ensuremath{\equiv}}}1
    {≃}{{\ensuremath{\simeq}}}1
    {≅}{{\ensuremath{\cong}}}1
    {×}{{\ensuremath{\times}}}1
    {·}{{\ensuremath{\cdot}}}1
    {±}{{\ensuremath{\pm}}}1
    {⊕}{{\ensuremath{\oplus}}}1
    {⊗}{{\ensuremath{\otimes}}}1
    {⋆}{{\ensuremath{\star}}}1
    {▸}{{\ensuremath{\blacktriangleright}}}1
    {⟨}{{\ensuremath{\langle}}}1
    {⟩}{{\ensuremath{\rangle}}}1
    {⦃}{{\ensuremath{\{\!\mid}}}1
    {⦄}{{\ensuremath{\mid\!\}}}}1
    {⟦}{{\ensuremath{[\![}}}1
    {⟧}{{\ensuremath{]\!]}}}1
    {⌊}{{\ensuremath{\lfloor}}}1
    {⌋}{{\ensuremath{\rfloor}}}1
    {⌈}{{\ensuremath{\lceil}}}1
    {⌉}{{\ensuremath{\rceil}}}1
    {√}{{\ensuremath{\surd}}}1
    {∣}{{\ensuremath{\mid}}}1
    {∤}{{\ensuremath{\nmid}}}1
    {∎}{{\ensuremath{\blacksquare}}}1
    {⋯}{{\ensuremath{\cdots}}}1
    {…}{{\ensuremath{\ldots}}}1
    {∑}{{\ensuremath{\sum}}}1
    {∏}{{\ensuremath{\prod}}}1
    {⁻¹}{{\ensuremath{^{-1}}}}1
    {¹}{{\ensuremath{^1}}}1
    {²}{{\ensuremath{^2}}}1
    {³}{{\ensuremath{^3}}}1
    {ⁿ}{{\ensuremath{^n}}}1
    {ᵐ}{{\ensuremath{^m}}}1
    {₀}{{\ensuremath{_0}}}1
    {₁}{{\ensuremath{_1}}}1
    {₂}{{\ensuremath{_2}}}1
    {₃}{{\ensuremath{_3}}}1
    {₄}{{\ensuremath{_4}}}1
    {₅}{{\ensuremath{_5}}}1
    {ₙ}{{\ensuremath{_n}}}1
    {ₘ}{{\ensuremath{_m}}}1
    {ₖ}{{\ensuremath{_k}}}1
    {ᵢ}{{\ensuremath{_i}}}1
    {ⱼ}{{\ensuremath{_j}}}1
    {ℤ}{{\ensuremath{\mathbb{Z}}}}1
    {ℕ}{{\ensuremath{\mathbb{N}}}}1
    {ℚ}{{\ensuremath{\mathbb{Q}}}}1
    {ℝ}{{\ensuremath{\mathbb{R}}}}1
    {ℂ}{{\ensuremath{\mathbb{C}}}}1
    {𝔽}{{\ensuremath{\mathbb{F}}}}1
    {𝒫}{{\ensuremath{\mathcal{P}}}}1
    {∗}{{\ensuremath{\ast}}}1
    {•}{{\ensuremath{\bullet}}}1,
  % Comments
  morecomment=[l]{--},
  morecomment=[s]{/-}{-/},
  morestring=[b]",
  stringstyle=\color{stringcolor},
  keepspaces=true,
  keywordstyle=[1]\color{keywordcolor}\bfseries,
  keywordstyle=[2]\color{sortcolor}\bfseries,
  keywordstyle=[3]\color{tacticcolor}\bfseries,
  keywordstyle=[4]\color{errorcolor}\bfseries\underline,
  escapeinside={(*@}{@*)},
  columns=fullflexible,
  basicstyle=\ttfamily\small,
  commentstyle=\color{commentcolor}\itshape,
  breaklines=true,
  showstringspaces=false,
  frame=single,
  framesep=2mm,
  xleftmargin=4mm,
  numbers=left,
  numberstyle=\tiny\color{commentcolor},
  numbersep=6pt,
}

% Inline Lean code in prose: \lean{Int.ModEq.pow}
\newcommand{\lean}[1]{\lstinline[language=lean,basicstyle=\ttfamily\normalsize]|#1|}

\lstset{language=lean}
% --- XeTeX rendering patch (verbatim from claims/tatami-mod8-defect/claim.tex
%     2026-09-27 sanctioned block, itself backported from
%     papers/cyl3-mod23/paper.tex; extended with the 6 further math symbols
%     of this frozen file's inventory: 00D7 2083 2099 2212 2223 27FA;
%     00B7/2014/2026 are in xeCJK's punctuation lists (MiddlePunct/LongPunct,
%     xeCJK.sty) and are left to xeCJK's own CJK font routing).
%     Rendering-only; source bytes of listings unchanged. Guarded so pdflatex
%     keeps literate-table behaviour.
\ifdefined\XeTeXversion
  {\lccode`\~="2022 \lowercase{\global\catcode`~=\active \global\def~{\ensuremath{\bullet}}}}
  {\lccode`\~="2115 \lowercase{\global\catcode`~=\active \global\def~{\ensuremath{\mathbb{N}}}}}
  {\lccode`\~="2124 \lowercase{\global\catcode`~=\active \global\def~{\ensuremath{\mathbb{Z}}}}}
  {\lccode`\~="2190 \lowercase{\global\catcode`~=\active \global\def~{\ensuremath{\leftarrow}}}}
  {\lccode`\~="2192 \lowercase{\global\catcode`~=\active \global\def~{\ensuremath{\rightarrow}}}}
  {\lccode`\~="2194 \lowercase{\global\catcode`~=\active \global\def~{\ensuremath{\leftrightarrow}}}}
  {\lccode`\~="2200 \lowercase{\global\catcode`~=\active \global\def~{\ensuremath{\forall}}}}
  {\lccode`\~="2203 \lowercase{\global\catcode`~=\active \global\def~{\ensuremath{\exists}}}}
  {\lccode`\~="2227 \lowercase{\global\catcode`~=\active \global\def~{\ensuremath{\wedge}}}}
  {\lccode`\~="2228 \lowercase{\global\catcode`~=\active \global\def~{\ensuremath{\vee}}}}
  {\lccode`\~="2261 \lowercase{\global\catcode`~=\active \global\def~{\ensuremath{\equiv}}}}
  {\lccode`\~="2264 \lowercase{\global\catcode`~=\active \global\def~{\ensuremath{\leq}}}}
  {\lccode`\~="2265 \lowercase{\global\catcode`~=\active \global\def~{\ensuremath{\geq}}}}
  {\lccode`\~="22A2 \lowercase{\global\catcode`~=\active \global\def~{\ensuremath{\vdash}}}}
  {\lccode`\~="27E8 \lowercase{\global\catcode`~=\active \global\def~{\ensuremath{\langle}}}}
  {\lccode`\~="27E9 \lowercase{\global\catcode`~=\active \global\def~{\ensuremath{\rangle}}}}
  {\lccode`\~="00D7 \lowercase{\global\catcode`~=\active \global\def~{\ensuremath{\times}}}}
  {\lccode`\~="2083 \lowercase{\global\catcode`~=\active \global\def~{\ensuremath{{}_3}}}}
  {\lccode`\~="2099 \lowercase{\global\catcode`~=\active \global\def~{\ensuremath{{}_n}}}}
  {\lccode`\~="2212 \lowercase{\global\catcode`~=\active \global\def~{\ensuremath{-}}}}
  {\lccode`\~="2223 \lowercase{\global\catcode`~=\active \global\def~{\ensuremath{\mid}}}}
  {\lccode`\~="27FA \lowercase{\global\catcode`~=\active \global\def~{\ensuremath{\Longleftrightarrow}}}}
\fi
\ifdefined\XeTeXversion
  \usepackage{xeCJK}
  \setCJKmainfont{SimSun}
\fi

\newtheorem{theorem}{Theorem}

\title{Mod-2 period-6 and mod-3 anti-period laws for the recurrence family
$(6,9,0,-1)$, with a matrix-functional generalization\\
(Lean 4 machine-checked claim of record)}
\author{Aurora0134\thanks{Contact: via the GitHub account Aurora0134; ORCID to
be added at Zenodo upload.}}
\date{September 27, 2026}

\begin{document}
\maketitle

\begin{abstract}
We record two theorems about the order-4 integer recurrence $c(n+4) =
6c(n+3) + 9c(n+2) - c(n)$ --- the recurrence of the matching counts of the
$3 \times n$ cylindrical grid (OEIS entry A033515). Theorem~1: every
integer sequence satisfying this recurrence has period dividing $6$ modulo
$2$, and satisfies the anti-period relation $c(n+4) \equiv -c(n) \pmod 3$.
Theorem~2: every matrix-power functional of an $8 \times 8$ integer matrix
annihilated by $x^4 - 6x^3 - 9x^2 + 1$ inherits the same recurrence,
uniformly in the vector and the coordinate. All proofs are machine-checked
by the Lean 4 kernel (0 \texttt{sorry}; axioms within \{propext, Quot.sound,
Classical.choice\}). A structured novelty search found no occupying record.
The artifact package is deposited on Zenodo. This note is a priority claim
of record, not a full paper.
\end{abstract}

\section{Introduction}

The graph $C_3 \mathbin{\square} P_n$ --- the $3 \times n$ grid wrapped
around a cylinder, also called the stacked prism graph --- has a classical
matching count, OEIS entry A033515 \cite{oeis-a033515}, whose record lists
the order-4 recurrence $c(n+4) = 6c(n+3) + 9c(n+2) - c(n)$. Numerical
probes of three fixed three-hole defect configurations of the same grid
found (computationally; see Section~\ref{sec:scope}) counting sequences
that satisfy the same recurrence with different initial values.

This note records two theorems about that recurrence family, produced by
AI4Math, a highly automated LLM-plus-kernel pipeline in which the Lean 4
kernel is the sole arbiter of correctness. The present text is a priority
claim of record --- a short deposit note accompanying a Zenodo artifact
package; a fuller narrative paper is prepared separately through the
pipeline's paper lane. Every theorem stated below is a \emph{complete
proof} in the pipeline's grading: it compiles with 0 \texttt{sorry}, its
axioms stay within the standard whitelist, and the statement was frozen
before proving.

\section{Statement}

The formal development quantifies over an \emph{arbitrary} integer sequence
satisfying the recurrence, so Theorem~\ref{thm:master} covers the whole
family (all initial values) at once; congruences use mathlib's
\texttt{Int.ModEq} notation \texttt{[ZMOD m]}. The listings below reproduce
the frozen theorem signatures verbatim (byte-identical excerpts of the
frozen statement files \texttt{04-statement-p1.lean} and
\texttt{04-statement-p2.lean}; the files' placeholder \texttt{sorry} bodies
are not reproduced; the Chinese docstrings omitted here remain in the
package's \texttt{proofs/} directory).

% LEAN: tasks/20260926-cyl3-pipeline :: cyl3_master grade=完全证明
\begin{theorem}[recurrence-family congruence law]\label{thm:master}
Let $c : \mathbb{N} \to \mathbb{Z}$ satisfy $c(n+4) = 6c(n+3) + 9c(n+2) -
c(n)$ for all $n$. Then, for all $n$,
\[
  c(n+6) \;\equiv\; c(n) \pmod 2
  \qquad\text{and}\qquad
  c(n+4) \;\equiv\; -c(n) \pmod 3.
\]
\end{theorem}

\begin{lstlisting}[caption={Formal statement of Theorem~\ref{thm:master} (verbatim excerpt, lines 11--13 of the frozen \texttt{formalized/04-statement-p1.lean}, sha256 \texttt{4740b9d5...e733bc}).},label={lst:p1}]
theorem cyl3_master (c : ℕ → ℤ)
    (h : ∀ n, c (n+4) = 6 * c (n+3) + 9 * c (n+2) - c n) :
    (∀ n, c (n+6) ≡ c n [ZMOD 2]) ∧ (∀ n, c (n+4) ≡ -c n [ZMOD 3]) := by
\end{lstlisting}

% LEAN: tasks/20260926-cyl3-pipeline :: cyl3_matrix_family grade=完全证明
\begin{theorem}[matrix-level inheritance]\label{thm:matrix}
Let $T$ be an $8 \times 8$ integer matrix with $T^4 = 6T^3 + 9T^2 - I$,
let $v$ be an integer vector of length $8$, and let $i$ be a coordinate.
Then the sequence $a(n) = ((T^n)\,v)_i$ satisfies the same recurrence: for
all $n$,
\[
  a(n+4) \;=\; 6\,a(n+3) \;+\; 9\,a(n+2) \;-\; a(n).
\]
\end{theorem}

\begin{lstlisting}[caption={Formal statement of Theorem~\ref{thm:matrix} (verbatim excerpt, lines 13--18 of the frozen \texttt{formalized/04-statement-p2.lean}, sha256 \texttt{8cac6fa4...c1d7fe}).},label={lst:p2}]
theorem cyl3_matrix_family {T : Matrix (Fin 8) (Fin 8) ℤ}
    (v : Fin 8 → ℤ) (i : Fin 8)
    (hT : T ^ 4 = 6 • T ^ 3 + 9 • T ^ 2 - (1 : Matrix (Fin 8) (Fin 8) ℤ))
    (n : ℕ) :
    (T ^ (n+4)).mulVec v i
      = 6 * (T ^ (n+3)).mulVec v i + 9 * (T ^ (n+2)).mulVec v i - (T ^ n).mulVec v i := by
\end{lstlisting}

\section{Proof}

The complete machine proofs are reproduced below verbatim:
Listing~\ref{lst:full} is the entire frozen final development
\texttt{final/04-proof-complete.lean} (69 lines, sha256
\texttt{566e9e1d...c5ae8}), mechanically spliced byte-for-byte. Its Chinese
header comments are pipeline provenance notes (route sketches and source
pointers), not mathematical content. A narrative prose proof is not the
duty of this note; it belongs to the full paper lane.

\begin{lstlisting}[caption={Complete development, verbatim: lines 1--69 (the whole file) of the frozen \texttt{final/04-proof-complete.lean}, sha256 \texttt{566e9e1d...c5ae8}.},label={lst:full}]
import Mathlib.Data.Int.ModEq
import Mathlib.Data.Matrix.Mul

/-!
# 04-proof-complete（军政部 r1）

两定理定理头与 tasks/20260926-cyl3-pipeline/formalized/04-statement-p{1,2}.lean
冻结文本逐字一致；本文件只写 proof body。AI 生成（dept-04-prove，
node=anthropic/a6api-main/kimi-k3），kernel 实编裁决为准。

- P1 路线（03 同型 r2 实证形态）：Int.modEq_iff_dvd 显式 dvd 证人 + ring；
  模 2 由 `c (n+4) ≡ c (n+2) + c n` 在 n 与 n+2 两处展开后 add_right/trans 收束；
  模 3 反周期一步证人（6、9 被 3 整除）。
- P2 路线（EdgemidSmoke03r3.lean edgemid_matrix_family 同型临摹，系数 6,9,−1）：
  hgen = pow_add + hT + Matrix.mul_sub/mul_add + mul_smul_comm×2 + Matrix.mul_one
  + ← pow_add×2；mulVec 线性 = Matrix.sub_mulVec/add_mulVec + smul_mulVec×2；
  收尾 first | rfl | simp only [smul_eq_mul] | …（03 r3 已编译验证）。
-/

/-- P1（`cyl3_master`，修正档=稳妥）：任一满足四阶线性递推
`c (n+4) = 6 * c (n+3) + 9 * c (n+2) - c n` 的整数序列，
模 2 以 6 为周期，且模 3 有反周期 `c (n+4) ≡ -c n [ZMOD 3]`
（推论：`3 ∣ c n ⟺ 3 ∣ c (n % 4)` 型整除刻画与基线奇偶模式）。
人话语义：带缺陷圆柱格 C₃×Pₙ（3×n，三种固定孔位）匹配计数族的
统一同余律——无孔基线（OEIS A033515）与同行相邻/对角/三角三种
三孔缺陷序列同满足递推 (6,9,0,−1)，本定理对该递推整族一次覆盖。
题源：tasks/20260926-beian-select/candidates/04-cylinder-3holes-review.md（收窄命题 P1）。 -/
theorem cyl3_master (c : ℕ → ℤ)
    (h : ∀ n, c (n+4) = 6 * c (n+3) + 9 * c (n+2) - c n) :
    (∀ n, c (n+6) ≡ c n [ZMOD 2]) ∧ (∀ n, c (n+4) ≡ -c n [ZMOD 3]) := by
  have key : ∀ m, c (m+4) ≡ c (m+2) + c m [ZMOD 2] := by
    intro m
    refine Int.modEq_iff_dvd.2 ⟨-(3 * c (m+3) + 4 * c (m+2) - c m), ?_⟩
    rw [h m]
    ring
  refine ⟨fun n => ?_, fun n => ?_⟩
  · have h1 : c (n+4) ≡ c (n+2) + c n [ZMOD 2] := key n
    have h2 : c (n+6) ≡ c (n+4) + c (n+2) [ZMOD 2] := key (n+2)
    have h3 : c (n+6) ≡ (c (n+2) + c n) + c (n+2) [ZMOD 2] := h2.trans (h1.add_right _)
    exact h3.trans (Int.modEq_iff_dvd.2 ⟨-c (n+2), by ring⟩)
  · refine Int.modEq_iff_dvd.2 ⟨-(2 * c (n+3) + 3 * c (n+2)), ?_⟩
    rw [h n]
    ring

/-- P2（`cyl3_matrix_family`，挑战档）：8×8 整数转移矩阵 `T` 一旦被
四次多项式零化（`T ^ 4 = 6 • T ^ 3 + 9 • T ^ 2 - 1`），其任意矩阵幂
泛函 `(T ^ (n+k)).mulVec v i` 都满足同一个四阶递推
`a (n+4) = 6 * a (n+3) + 9 * a (n+2) - a n`——一条引理统一三种
固定孔位缺陷计数在转移矩阵层的共有结构（矩阵泛函层统一递推定理）。
人话语义：这是转移矩阵法（Lundow 1996/1998 口径）的泛函形式化；
「该矩阵确实等于缺陷匹配计数」的组合语义桥不进 kernel，
论文/报告层按主张降级如实声明，不得拔高为计数本体定理。
题源：tasks/20260926-beian-select/candidates/04-cylinder-3holes-review.md（收窄命题 P2）。 -/
theorem cyl3_matrix_family {T : Matrix (Fin 8) (Fin 8) ℤ}
    (v : Fin 8 → ℤ) (i : Fin 8)
    (hT : T ^ 4 = 6 • T ^ 3 + 9 • T ^ 2 - (1 : Matrix (Fin 8) (Fin 8) ℤ))
    (n : ℕ) :
    (T ^ (n+4)).mulVec v i
      = 6 * (T ^ (n+3)).mulVec v i + 9 * (T ^ (n+2)).mulVec v i - (T ^ n).mulVec v i := by
  have hgen : T ^ (n+4)
      = 6 • T ^ (n+3) + 9 • T ^ (n+2) - (T^n : Matrix (Fin 8) (Fin 8) ℤ) := by
    rw [pow_add, hT, Matrix.mul_sub, Matrix.mul_add,
        mul_smul_comm, mul_smul_comm,
        Matrix.mul_one, ← pow_add, ← pow_add]
  -- 收尾：rw 自带 rfl 不足以穿透 Pi 逐点作用（r2 实测 unsolved goals），
  -- 显式 `rfl` tactic 以全 defeq 收官（r1 级联实测首支即中，故仅留此一支）。
  rw [hgen, Matrix.sub_mulVec, Matrix.add_mulVec,
      Matrix.smul_mulVec, Matrix.smul_mulVec]
  rfl
\end{lstlisting}

\section{Verification evidence}

\begin{itemize}
\item Toolchain: Lean v4.34.0 (\texttt{leanprover/lean4:v4.34.0})
\cite{lean4}, mathlib4 v4.34.0 (rev \texttt{5ed29652}) \cite{mathlib}.
\item Statement integrity: frozen statements
\texttt{formalized/04-statement-p1.lean}, sha256
\texttt{4740b9d53e97360c49d0eefe97a7abc470134acf8e61273963d7a55de8e733bc},
and \texttt{formalized/04-statement-p2.lean}, sha256
\texttt{8cac6fa4ec068dc913ff90769bda34333bddd9d22c7463d2cecab0e3edc1d7fe}
(LF-normalized); final development
\texttt{final/04-proof-complete.lean}, sha256
\texttt{566e9e1d7053602eda7853943f582e3d3726009c32f9dc4b6ce3ca8ee19c5ae8};
symmetric difference of the two theorem heads (up to \texttt{:= by})
against the frozen basis: zero (verbatim).
\item Kernel check: strict re-compilation (\texttt{scripts/lean-verify},
strict mode) exit 0 with empty diagnostics; a byte-level scan of the whole
file (comments included) finds 0 occurrences of
\texttt{sorry}/\texttt{admit}. Axiom audit (sidecar re-compilation),
verbatim output (\texttt{audit/audit-axioms-sidecar.log}):
\begin{verbatim}
'cyl3_master' depends on axioms: [propext, Classical.choice, Quot.sound]
'cyl3_matrix_family' depends on axioms: [propext, Classical.choice, Quot.sound]
\end{verbatim}
Both are within the standard whitelist \{propext, Quot.sound,
Classical.choice\}; no \texttt{sorryAx}.
\item Degeneracy probing: the well-definedness gate ran 14 bare-tactic
probes (7 against each statement: \texttt{nlinarith}, \texttt{norm\_num},
\texttt{simp}, \texttt{ring\_nf}, \texttt{positivity}, \texttt{omega},
\texttt{linarith}); all failed as required, so neither theorem is closable
by a bare decision procedure.
\item Audit chain: gate-2 well-definedness compilation plus degeneracy
probing, statement freeze with sha256 pins, gate-3 final strict compilation
and axiom audit (\texttt{audit/final-04-audit.md}: all four checks pass;
grading \emph{complete proof} for both theorems). Evidence originals are in
the artifact package \texttt{audit/}.
\end{itemize}

\section{Novelty evidence}

The pipeline's C9 novelty re-review (2026-09-26; independent re-query
before deposit) adjudicated, per proposition: (A) the three fixed
three-hole defect sequences --- \emph{no occupying record found},
established and strengthened (the sequences were recomputed by an
independent enumerator implementation, then probed on OEIS through 19
sequence-shape queries and 6 keyword groups and on the literature channels
below); (P1) Theorem~\ref{thm:master} --- \emph{no occupying record found},
positioned as a congruence-law observation accompanying the new sequences
(new-sequence tier; an elementary Pisano-type instance of classical
modular-period theory, not new mathematics); (P2)
Theorem~\ref{thm:matrix} --- \emph{no occupying record found}, positioned
as a library-lemma-level first formalization; (P3) the full combinatorial
semantics --- occupied at the literature layer
\cite{lundow1996,lundow1998,oeis-a033515,oeis-a287428}, hence not a
contribution of this note. This is search evidence, not a claim to new
mathematics. Query-channel summary (the full query log, including zero-hit
evidence files, is copied verbatim into \texttt{audit/c9-record.md} of this
package):

\begin{center}
\begin{tabular}{p{0.24\textwidth}p{0.46\textwidth}p{0.20\textwidth}}
\hline
channel & query summary & result \\
\hline
OEIS sequence layer & 26 queries (19 sequence-shape probes of the defect
families plus recurrence-text and generating-function probes) + 2 record
reads; calibration query hits A033515 (channel works) & defect variants:
all null \\
arXiv & 8 query groups (incl.\ 3 http trial-and-error); nearest lines
arXiv:2109.12716 \cite{dey-krishnan}, cond-mat/0611449 \cite{kong},
2406.05750 \cite{arora} read and delineated & no occupant \\
OpenAlex & 5 broad + 3 title-exact queries & no occupant \\
zbMATH & 4 query groups & no occupant; 3 records license-masked (gap) \\
Semantic Scholar & 1 query & HTTP 429, not retried (gap) \\
general web (Bing/DDG) & 5 fetches & channel degraded (gap) \\
CNKI & not searched (login required) & gap \\
formal libraries & mathlib / compfiles / sequencelib probes
(\texttt{Hosoya}, \texttt{monomer}, \texttt{A033515}, \dots) & zero hits;
A033515 absent from sequencelib \\
\hline
\end{tabular}
\end{center}

Recorded residual gaps, verbatim in substance from the re-review: CNKI not
searched (login required); one Semantic Scholar query rate-limited (HTTP
429) and not retried; three zbMATH records license-masked and unreadable;
the general web search layer degraded (Bing localization pollution;
DuckDuckGo blocked by DNS pollution and proxy failure).

\section{Scope and limitations}\label{sec:scope}

Grading boundary, per the final audit's transfer wording: the two theorems
are \emph{complete proofs}, strictly as statements about the recurrence
family (an arbitrary integer sequence $c$) and about matrix-power
functionals. The kernel did \emph{not} prove that any concrete defect
matching-count sequence equals this recurrence or this matrix: the
statement ``the three fixed three-hole defect counting sequences satisfy
the recurrence $(6,9,0,-1)$'' rests on the OEIS record and the
transfer-matrix literature (A033515 and A287428 are registered, not new
\cite{oeis-a033515,oeis-a287428,lundow1996,lundow1998}) plus a local
computational probe (range $n \leq 14$, two independently implemented
enumerators agreeing term by term); it is reported at the
computational-conjecture layer and is not part of the kernel-proved
content. Probe values, copied verbatim from the source records:

\begin{center}
\begin{tabular}{lll}
\hline
family (fixed three-hole configurations) & first values & start \\
\hline
baseline, no holes (OEIS A033515, registered) &
$1, 4, 32, 228, 1655, 11978, 86731, 627960$ & $n \geq 0$ \\
row3: three adjacent holes in one row &
$22, 135, 1006, 7251, 52538, 380352, 2753948, 19939605$ & $n \geq 3$ \\
diag3: three diagonal holes &
$13, 83, 615, 4437, 32144, 232714, 1684965, 12199779$ & $n \geq 3$ \\
tri3: three triangular holes &
$3, 19, 141, 1017, 7368, 53342, 386223, 2796399$ & $n \geq 2$ \\
\hline
\end{tabular}
\end{center}

Further boundaries. Theorem~\ref{thm:matrix} is an elementary consequence
of the annihilation hypothesis (Cayley--Hamilton-type reasoning); the
contribution is the packaging as one reusable kernel-checked lemma ---
a library-lemma-level first formalization, not new mathematics. A fourth
configuration with a moving hole was recorded with a sixth-order recurrence
fit; it stays at record level and is not a contribution claim. The novelty
evidence above is negative search evidence with the stated residual gaps;
it supports no first-discovery claim beyond the stated tier. The source
pipeline report is marked ``pending gate-4 human sign-off''; this deposit
is issued under the depositing author's explicit start instruction of
2026-09-27 (recorded in this package's card). This note contains no
literature survey; it is a deposit record, not a survey article.

\section*{Data availability}

Artifact package: Zenodo deposit (DOI to be reserved at upload; see
UPLOAD.md in this package), with a public mirror on GitHub at
\url{https://github.com/Aurora0134/ai4math-claims}. Note under CC BY 4.0,
code under MIT.

\section*{Use of generative AI}

(a) Stage mapping: candidate-proposition generation, formalization into
Lean statements, proof search, and text drafting (including this note) were
performed by AI (LLM) stages of the AI4Math pipeline; topic adjudication,
statement-semantics verification and gate sign-off were human. (b)
Model/node/budget (numbers copied verbatim from the source task ledger
\texttt{tasks/20260926-cyl3-pipeline/budget.log}): a single resolved model
node, wire-id \texttt{anthropic/a6api-main/kimi-k3} (resolution level L2
relay), no cross-node switching; the source task totals llm-call $1/6$ (the
roundtrip semantic judge), kernel compilations $12/24$, diagnostic probes
$0/4$; the separate full-paper lane used LaTeX compilations $6/10$; the C9
exclusivity re-review used 0 LLM calls and 60 retrieval requests. This
claim-note stage used llm-call $0$ of its budget $4$, and $5$ LaTeX
compilations of its budget $6$ (counting the final rebuild that embedded
this number; ledger \texttt{claims/cyl3-mod23/budget.log}). (c) Final
arbiter: all statements and proofs reported here were checked by the Lean 4
kernel: every proof compiles with no \texttt{sorry}, and \texttt{\#print
axioms} reports only \{propext, Quot.sound, Classical.choice\} for both
theorems. (d) The human author reviewed and is responsible for all content;
the AI system is not an author.

\section*{Author contributions}

CRediT-style: the human author was responsible for topic selection and
adjudication, verification of statement semantics, gate sign-off, and the
decision to deposit; the AI4Math pipeline (an LLM system) performed
proposition generation, proof search, and text drafting. The AI system is
not listed as an author.

\begin{thebibliography}{9}

\bibitem{lean4}
L.~de Moura and S.~Ullrich.
\newblock The Lean 4 theorem prover and programming language.
\newblock \emph{CADE-28}, LNCS 12699, pp.~625--635, 2021.
\url{https://doi.org/10.1007/978-3-030-79876-5_37}

\bibitem{mathlib}
The mathlib Community.
\newblock The Lean mathematical library.
\newblock \emph{CPP 2020}, pp.~367--381, 2020.
\url{https://doi.org/10.1145/3372885.3373824}

\bibitem{dey-krishnan}
P.~S. Dey and K.~Krishnan.
\newblock Disordered monomer-dimer model on cylinder graphs.
\newblock \emph{arXiv:2109.12716 [math.PR]}, 2021.
\url{https://arxiv.org/abs/2109.12716}

\bibitem{kong}
Y.~Kong.
\newblock Exact asymptotics of monomer-dimer model on rectangular
semi-infinite lattices.
\newblock \emph{arXiv:cond-mat/0611449 [cond-mat.stat-mech]}, 2006.
\url{https://arxiv.org/abs/cond-mat/0611449}

\bibitem{arora}
A.~Arora.
\newblock The monopole-dimer model on high-dimensional cylindrical, toroidal,
M\"obius and Klein grids.
\newblock \emph{arXiv:2406.05750 [math.CO]}, 2024.
\url{https://arxiv.org/abs/2406.05750}

\bibitem{lundow1996}
P.~H. Lundow.
\newblock Computation of matching polynomials and the number of 1-factors in
polygraphs.
\newblock Research reports No.~12, Department of Mathematics, Ume{\aa}
University, 1996; as referenced in OEIS A033515 \cite{oeis-a033515}.

\bibitem{lundow1998}
P.~H. Lundow.
\newblock Enumeration of matchings in polygraphs.
\newblock 1998; as referenced in OEIS A033515 \cite{oeis-a033515}.

\bibitem{oeis-a033515}
The OEIS Foundation.
\newblock Entry A033515 (P.~H. Lundow): number of matchings in graph
$C_3 \times P_n$.
\newblock \emph{The On-Line Encyclopedia of Integer Sequences}.
\url{https://oeis.org/A033515}

\bibitem{oeis-a287428}
The OEIS Foundation.
\newblock Entry A287428 (A.~Howroyd, 2017): $T(m,n)$, number of matchings in
the stacked prism graph $C_m \times P_n$.
\newblock \emph{The On-Line Encyclopedia of Integer Sequences}.
\url{https://oeis.org/A287428}

\end{thebibliography}

\end{document}
